\chapter{Perception}
\label{ch:perception}

\pkg{dharma_perception} turns the LiDAR into two sector scans in \param{robot.frames.base}: a
front scan for reactive avoidance (\topic{perception/sector_scan}) and a fail-safe 360° scan for
the protective field and the follower's motion guard (\topic{perception/safety_scan}). It also
throttles the camera for the fiducial detector.

\section{Sector scan semantics}

A \code{SectorScan} holds the nearest obstacle per angular sector; sector $i$ is centred at
\code{angle_min} $+ (i + \tfrac12)$ \code{angle_increment}. \texttt{ranges[i] >= range\_max} means
nothing within range; \texttt{ranges[i] <= BLOCKED} (0) or NaN means blocked or unknown, and
\emph{every consumer treats it as an obstacle at the robot}. Every sector is \code{BLOCKED}
when:
\begin{itemize}
  \item the LiDAR is blind: fewer than \param{blind.min_points} points off the robot body, or
    (with \param{blind.min_region_points}) one of \param{blind.regions} azimuth regions empty;
  \item the sensor mount transform or, with \param{filter.level_source} \code{tf}, the roll and
    pitch of the base (\frm{odom}$\to$\frm{base_link}) is missing;
  \item processing raised an exception, or no usable cloud arrived for
    \param{input.timeout_sec} (driver down, only malformed messages).
\end{itemize}
In 0.1 an all-zero scan was read as free by the protective field (review B1) and NaN sectors as
free by the follower (F1); this convention removes both.

\section{Pipeline}

\begin{enumerate}
  \item \textbf{Parse.} The serialized Livox \code{CustomMsg} is read with numpy (no per-point
    Python objects; every length validated, both endiannesses). PointCloud2 goes through
    \code{dharma_utils.ros.pointcloud.read_xyz}. A malformed message is dropped with a
    throttled error (review \#4).
  \item \textbf{Livox noise.} The tag byte marks each point per field (bits 0--1 adhesion
    between objects, 2--3 rain, fog or dust, 4--5 other; 0 high, 1 medium, 2 low confidence).
    Points at \param{filter.noise_level} in the enabled fields are dropped (review \#12).
  \item \textbf{Self filter.} Points inside \param{robot.footprint} grown by
    \param{filter.self_margin}.
  \item \textbf{Level.} Points in \frm{base_link} are rotated by the roll and pitch of the base
    in \frm{odom} at the cloud stamp, so the band is relative to the floor while the gait pitches
    the body (finding P1).
  \item \textbf{Height band.} \param{filter.height_min} to \param{filter.height_max} above the
    floor, the floor being \param{robot.base_height} below \frm{base_link}. Optional
    \param{filter.ground_segmentation} estimates the floor per cell (lowest point within a
    maximum step) for ramps and thresholds.
  \item \textbf{Sectors.} Each range is the $k$-th nearest point of the sector (\param{front.k},
    \param{safety.k}, default 2): a single dust return does not stop the robot; a sector with
    fewer than $k$ points is free.
  \item \textbf{Obstacle memory.} Confirmed safety points are kept in \frm{odom} for
    \param{memory.memory_sec} and merged into the safety scan while they are hidden below the
    lowest beam; an odometry jump clears the memory (review B4).
\end{enumerate}

\textbf{Why levelling matters.} For a flat floor and \frm{base_link} at height $h$ pitched by
$\vartheta$, a floor point at horizontal range $r$ appears at $z \approx -h + r\sin\vartheta$ and
enters the band at
\begin{equation}
  r > r_{\mathrm{false}} = \frac{h + z_{\min}}{\sin\vartheta}.
\end{equation}
\begin{center}
\begin{tabular}{@{}lrrrr@{}}
  \toprule
  Band margin $h + z_{\min}$ & $\vartheta = 1^\circ$ & $2^\circ$ & $3^\circ$ & $5^\circ$ \\
  \midrule
  10~cm & 5.7~m & 2.9~m & 1.9~m & 1.15~m \\
  5~cm & 2.9~m & 1.4~m & 0.96~m & 0.57~m \\
  \bottomrule
\end{tabular}
\end{center}
A trotting legged robot pitches by a few degrees, so in \frm{base_link} the floor became an
obstacle within the 4~m safety range (false stops), and pitching up hid low obstacles (missed
detections). The 0.1 configuration also assumed $h = 0.35$~m in one file and 0.30~m in another;
both now come from the profile.

\section{Limits}
\label{sec:perception-limits}

\begin{itemize}
  \item \textbf{Near-field blind zone.} A Mid-360 sees down to $-7^\circ$: mounted 0.53~m above
    the floor it sees a 10~cm object only beyond about 3.5~m. The memory keeps objects seen
    earlier, but an object that appears inside the zone (placed there, or the robot switched on
    next to it) is not seen. Use a pitched mount or a near-field sensor (bumper, ToF, 2D
    scanner) where this matters.
  \item \textbf{Negative obstacles} (stairs down, holes, platform edges) are not detected:
    missing floor is not an obstacle.
  \item \textbf{Motion distortion.} The scan is not deskewed: a Mid-360 frame spans 100~ms, so
    a point is up to $v \cdot 0.1$~m off (8~cm at 0.8~m/s) and turning at $\omega$ moves a point
    at range $r$ by up to $0.1\,\omega r$. The latency from the first point to the published scan
    is the frame (100~ms), driver and transport (about 5--10~ms) and processing (about
    10--15~ms for 20k points on one desktop core); \param{robot.limits.reaction_time} must stay
    above it and the field margins above the distortion.
  \item \textbf{Quantization.} The nearest point of a sector stands for the whole sector: 72
    safety sectors are 5° wide (lateral error up to $r\Delta\theta/2$: 8.7~cm at 2~m, 17.5~cm at
    4~m). Raise \param{safety.num_sectors} for a finer field.
  \item This is a software channel; it does not replace a safety-rated scanner
    (Section~\ref{sec:two-channel}).
\end{itemize}

\section{Nodes and parameters}

\begin{center}
\begin{tabularx}{\textwidth}{@{}p{0.30\textwidth}X@{}}
  \toprule
  Node & Role \\
  \midrule
  \code{lidar_sector_node} & Livox \code{CustomMsg} (raw bytes, default) or PointCloud2 in;
    \topic{perception/sector_scan}, \topic{perception/safety_scan}, \topic{perception/obstacles}
    out \\
  \code{apriltag_throttle} & camera image and info in, image at \code{rate_hz} out \\
  \code{apriltag} (\code{apriltag_ros}) & tag frames on TF for localization \\
  \bottomrule
\end{tabularx}
\end{center}

\readonlynote

% Generated by docs/manual/tools/reference.py; do not edit.
\begin{small}
\begin{longtable}{@{}p{0.28\textwidth}p{0.12\textwidth}p{0.08\textwidth}p{0.11\textwidth}p{0.29\textwidth}@{}}
\toprule
Parameter & Default & Unit & Range & Description \\
\midrule
\endhead
\bottomrule
\endfoot
\multicolumn{5}{@{}l}{\textbf{lidar\_sector\_node}} \\
\code{input.type}$^\dagger$ & \code{livox} &  & livox, pointcloud2 & livox: raw Livox CustomMsg bytes; pointcloud2: any PointCloud2 \\
\code{input.livox_topic}$^\dagger$ & \code{livox/lidar} &  &  & Livox CustomMsg topic \\
\code{input.pointcloud_topic}$^\dagger$ & \code{/cloud_registered_body} &  &  & PointCloud2 topic (frame fixed to the robot) \\
\code{input.frame}$^\dagger$ & \code{livox_frame} &  &  & sensor frame when the header has none \\
\code{input.timeout_sec}$^\dagger$ & 0.5 & s & 0.05\,..\,10 & no usable cloud for this long: every sector is published BLOCKED \\
\code{front.num_sectors}$^\dagger$ & 32 &  & 1\,..\,1440 & sectors over the fov \\
\code{front.fov_deg}$^\dagger$ & 180 & deg & 1\,..\,360 & centred ahead \\
\code{front.range_min}$^\dagger$ & 0.1 & m & 0\,..\,10 & closer: ignored \\
\code{front.range_max}$^\dagger$ & 8 & m & 0.1\,..\,100 & free beyond \\
\code{front.k}$^\dagger$ & 2 &  & 1\,..\,100 & a sector needs k points; its range is the k-th nearest \\
\code{safety.enabled}$^\dagger$ & true &  &  & publish perception/safety\_scan \\
\code{safety.num_sectors}$^\dagger$ & 72 &  & 4\,..\,1440 & sectors over 360 deg \\
\code{safety.range_min}$^\dagger$ & 0.1 & m & 0\,..\,10 & closer: ignored \\
\code{safety.range_max}$^\dagger$ & 4 & m & 0.1\,..\,100 & free beyond \\
\code{safety.k}$^\dagger$ & 2 &  & 1\,..\,100 & a sector needs k points; its range is the k-th nearest \\
\code{filter.self_margin}$^\dagger$ & 0.05 & m & 0\,..\,1 & footprint grown by this is the robot \\
\code{filter.level_source}$^\dagger$ & \code{tf} &  & tf, none & tf: roll/pitch of base in odom (unknown -> BLOCKED); none: base assumed level \\
\code{filter.tilt_max_age}$^\dagger$ & 0.2 & s & 0.01\,..\,5 & oldest odom->base usable for leveling \\
\code{filter.height_min}$^\dagger$ & 0.1 & m & -1\,..\,3 & band bottom above the floor \\
\code{filter.height_max}$^\dagger$ & 1.6 & m & 0\,..\,5 & band top above the floor \\
\code{filter.ground_segmentation}$^\dagger$ & false &  &  & estimate the floor height per cell (ramps, uneven floors) \\
\code{filter.ground_cell}$^\dagger$ & 0.3 & m & 0.05\,..\,5 & ground segmentation cell size \\
\code{filter.ground_max_step}$^\dagger$ & 0.05 & m & 0\,..\,0.5 & largest local floor offset accepted \\
\code{filter.drop_adhesion_noise}$^\dagger$ & true &  &  & Livox tag bits 0-1 (glue points) \\
\code{filter.drop_particle_noise}$^\dagger$ & true &  &  & Livox tag bits 2-3 (rain, dust) \\
\code{filter.drop_other_noise}$^\dagger$ & false &  &  & Livox tag bits 4-5 (other) \\
\code{filter.noise_level}$^\dagger$ & 2 &  & 1\,..\,2 & drop at this tag confidence: 1 medium, 2 low \\
\code{memory.enabled}$^\dagger$ & true &  &  & remember obstacles that left the view \\
\code{memory.memory_sec}$^\dagger$ & 10 & s & 0\,..\,60 & remembered points decay after this \\
\code{memory.radius}$^\dagger$ & 5 & m & 0.1\,..\,20 & points this close are remembered (merged only while below the lowest beam) \\
\code{memory.resolution}$^\dagger$ & 0.05 & m & 0.01\,..\,1 & memory cell size \\
\code{memory.max_cells}$^\dagger$ & 20000 &  & 100\,..\,1e+06 & memory size bound \\
\code{memory.lower_fov_deg}$^\dagger$ & -7 & deg & -90\,..\,0 & lowest beam elevation of the sensor (Mid-360: -7), for the near-field blind zone \\
\code{memory.max_jump}$^\dagger$ & 0.3 & m & 0\,..\,10 & odom jump beyond the motion: forget \\
\code{blind.min_points}$^\dagger$ & 100 &  & 0\,..\,1e+06 & fewer valid points in a cloud: blind \\
\code{blind.regions}$^\dagger$ & 4 &  & 1\,..\,36 & azimuth regions around the sensor z axis \\
\code{blind.min_region_points}$^\dagger$ & 0 &  & 0\,..\,100000 & fewer points in any region: blind (0: off) \\
\code{publish_markers}$^\dagger$ & true &  &  & perception/obstacles for RViz \\
\end{longtable}
\end{small}

